% =====================================================
% IDIOMA Y CODIFICACIÓN
% =====================================================

\usepackage[spanish]{babel}

\usepackage[utf8]{inputenc}

\usepackage[T1]{fontenc}

% =====================================================
% IMAGEN OVERLAY
% =====================================================


\usepackage{tikz}


% =====================================================
% MATEMÁTICAS
% =====================================================

\usepackage{amsmath}


% =====================================================
% GRÁFICOS Y TABLAS
% =====================================================

\usepackage{graphicx}

\usepackage{booktabs}


% =====================================================
% CITAS Y BIBLIOGRAFÍA APA
% =====================================================

\usepackage{csquotes}

\usepackage[
backend=biber,
style=apa,
sorting=nyt
]{biblatex}

\DeclareLanguageMapping{spanish}{spanish-apa}

\addbibresource{05_bibliography/references.bib}


% =====================================================
% CÓDIGO FUENTE
% =====================================================

\usepackage{listings}
\usepackage{xcolor}

% Cambiar nombres del entorno listings
\renewcommand{\lstlistingname}{Código}
\renewcommand{\lstlistlistingname}{Lista de códigos}

% =====================================================
% ALGORITMOS
% =====================================================

\usepackage{algorithm}
\usepackage{algpseudocode}

% Cambiar nombres del entorno algoritmos
\floatname{algorithm}{Algoritmo}
\renewcommand{\listalgorithmname}{Lista de algoritmos}


% =====================================================
% FIGURAS
% =====================================================

\usepackage{caption}

\captionsetup{
    font=small,
    labelfont=bf,
    labelsep=period
}


% =====================================================
% ACRÓNIMOS
% =====================================================

\usepackage{acro}

\acsetup{
    first-style = long-short,
    list/name = Lista de Acrónimos
}


% =====================================================
% HIPERVÍNCULOS
% =====================================================

\usepackage{hyperref}

\hypersetup{
    pdftitle={\documenttitle},
    pdfauthor={\documentauthor},
    pdfsubject={Trabajo Fin de Grado},
    pdfkeywords={TFG, investigación},
    colorlinks=true,
    linkcolor=black,
    citecolor=black,
    urlcolor=blue
}


% =====================================================
% REFERENCIAS INTELIGENTES
% =====================================================

\usepackage{cleveref}